\section*{Chapitre \ref{ch:dv:variance-swaps-and-volatility-derivatives} --- Swaps de variance et produits dérivés de volatilité}

\begin{solution}{exo:dv:variance-swaps-and-volatility-derivatives:1}
$N_{\mathrm{var}}=50\,000/(2\times20)=1\,250$ par point de variance~; à 25~\% le swap verse
$1\,250\times(625-400)=281\,250$.
\end{solution}

\begin{solution}{exo:dv:variance-swaps-and-volatility-derivatives:2}
Le strip verse $-2\ln(S_T/F)$, qui n'égale la variance réalisée qu'après ajout de $2\int dS/S$~: les gains de la détention de
$2/S_t$ dans l'indice. Sans la transaction quotidienne, la position est un \omterm{def:dv:variance-swaps-and-volatility-derivatives:log}{contrat log}, dont le paiement dépend de la
valeur finale de l'indice, et non du chemin suivi.
\end{solution}

\begin{solution}{exo:dv:variance-swaps-and-volatility-derivatives:3}
$f'(S)=\frac2{S_0}\ln(S/S_0)$ et $f^{\prime\prime}(S)=\frac2{S_0S}$. Ses pondérations de strip valent $2/(S_0K)$, contre $2/K^2$ pour le
\omterm{def:dv:variance-swaps-and-volatility-derivatives:vs}{swap de variance}~: par rapport au \omterm{def:dv:variance-swaps-and-volatility-derivatives:vs}{swap de variance}, il pondère chaque strike par $K/S_0$, moins sous le spot et davantage au-dessus.
Sur un indice à skew vers la baisse, les puts chers comptent pour moins, si bien que le strike est plus bas.
\end{solution}

\begin{solution}{exo:dv:variance-swaps-and-volatility-derivatives:4}
Variance finale espérée $0.75\times18^2+0.25\times16^2=243+64=307$ contre $22^2=484$~: la valeur vaut $-177$ par unité de
\omterm{def:dv:variance-swaps-and-volatility-derivatives:varnot}{notionnel de variance}.
\end{solution}

\begin{solution}{exo:dv:variance-swaps-and-volatility-derivatives:5}
$V=0.04$, $\Var=0.0004$~: $0.2-0.0004/(8\times0.04^{1.5})=0.2-0.00625=19.4\%$.
\end{solution}

\begin{solution}{exo:dv:variance-swaps-and-volatility-derivatives:6}
L'ajustement est proportionnel à la variance de la variance réalisée, divisée par son niveau. Aux échéances courtes
la variance n'a pas le temps de bouger~; aux échéances longues, le retour à la moyenne ($\kappa=2.05$, une demi-vie d'environ quatre mois)
en fait disparaître ses mouvements par moyennage. La variance de cette moyenne est maximale lorsque l'échéance dure quelques demi-vies.
\end{solution}

\begin{solution}{exo:dv:variance-swaps-and-volatility-derivatives:7}
Le strip complet donne 16.32~\%~; les options valant au moins 0.05 vont de 87 à 106, et le chiffre tombe à 16.10~\%,
0.22 point de moins~: les ailes manquantes portaient une part de la variance.
\end{solution}

\begin{solution}{exo:dv:variance-swaps-and-volatility-derivatives:8}
La couverture n'est exacte que pour des trajectoires continues dans la plage du strip. Chaque saut laisse une perte de l'ordre de
$|r|^3/3$ de variance, et un krach traversant le strike le plus bas laisse le paiement croître alors que la couverture s'arrête. La
prime au-dessus du strip est le prix de ce risque de saut et de troncature (0.18 point en moyenne selon le modèle de
Merton du chapitre 13, avec une perte de 1.9 points sur trois mois dans un cas sur cent), et non un profit.
\end{solution}

\begin{solution}{pb:dv:variance-swaps-and-volatility-derivatives:1}
\textbf{1.} 22.1~\% et 19.2~\%.
\textbf{2.} Les puts~: 72~\% de la variance~; les strikes inférieurs à 80 à eux seuls 24~\%.
\textbf{3.} 2\,262 par point de variance~; 931\,000.
\textbf{4.} 1.6 à 1.8 points, selon l'espacement.
\textbf{5.} 20.8~\%.
\textbf{6.} 21.4~\%~: le modèle calibré manque les ailes de la surface, que le strip pondère fortement.
\textbf{7.} 19.8~\%.
\textbf{8.} 0.72, 1.62 et 1.10 points.
\textbf{9.} La volatilité de la variance $\eta$~: 1.24 points pour $\eta-0.1$, 2.01 pour $\eta+0.1$.
\textbf{10.} Son paiement est la racine carrée de la variance réalisée, fonction non linéaire de la variation quadratique
de la trajectoire~; seules les fonctions du prix terminal sont des strips statiques, et la variance est la seule fonctionnelle
de chemin qui se réduit à un tel strip.
\textbf{11.} 17.1, 18.1 et 20.5.
\textbf{12.} $20.4-18.1=2.3$ points.
\textbf{13.} Elle descend, de 110~\% à 80~\% du contrat à terme jusqu'à 94~\% à 175~\%.
\textbf{14.} Avec des \omterm{def:dv:variance-swaps-and-volatility-derivatives:vs}{swaps de variance}, en dynamique~: la sensibilité d'un \omterm{def:dv:variance-swaps-and-volatility-derivatives:volswap}{swap de volatilité} à la variance vaut environ
$1/(2\sigma)$, si bien que le desk détient des \omterm{def:dv:variance-swaps-and-volatility-derivatives:vs}{swaps de variance} dans ce ratio et l'ajuste au gré des mouvements de volatilité, une couverture dont
le coût est la \omterm{def:dv:stochastic-volatility:sv}{volatilité de la volatilité} que l'ajustement valorise.
\textbf{15.} Au minimum la correction de saut, 0.18 point dans le modèle à sauts du chapitre 13, plus une réserve pour les krachs
traversant le strip.
\textbf{16.} La volatilité est la racine carrée de la variance, et la moyenne d'une racine carrée est inférieure à la racine carrée de
la moyenne. L'écart, ici 1.6 points, est le prix de la \omterm{def:dv:stochastic-volatility:sv}{volatilité de la volatilité}.
\textbf{17.} La distribution de la variance réalisée, surtout la \omterm{def:dv:stochastic-volatility:sv}{volatilité de la volatilité} $\eta$.
\textbf{18.} Une réserve pour le modèle~: l'étendue de l'ajustement sur les valeurs plausibles de $\eta$ (environ 0.4 point dans chaque sens pour $\pm0.1$)
et le coût de la recouverture lors d'un pic de volatilité.
\textbf{19.} 1.62 points de volatilité~: 21.4~\% contre 19.8~\% à un an.
\textbf{20.} La variance se réplique par un strip statique et un delta~; la volatilité exige un modèle.
\end{solution}

\begin{solution}{iq:dv:variance-swaps-and-volatility-derivatives:1}
La formule d'Itô donne $\int\sigma^2dt=2\int dS/S-2\ln(S_T/S_0)$. Détenir un strip statique d'options hors de la monnaie
avec des pondérations $2/K^2$ (le paiement log), et échanger chaque jour $2/S_t$ du sous-jacent~: le strike est le prix du strip
par unité de temps.
\iqlookfor{le \omterm{def:dv:variance-swaps-and-volatility-derivatives:log}{contrat log}, les pondérations en $1/K^2$ et la jambe dynamique.}
\end{solution}

\begin{solution}{iq:dv:variance-swaps-and-volatility-derivatives:2}
Le strip pondère les options par $1/K^2$, si bien que les bas strikes comptent davantage, et sur un indice les bas strikes se traitent à des
volatilités implicites plus élevées. Le strike est une moyenne de la variance sur tout le smile, penchée vers la baisse.
\iqlookfor{les pondérations et le skew.}
\end{solution}

\begin{solution}{iq:dv:variance-swaps-and-volatility-derivatives:3}
En dessous, par l'inégalité de Jensen. Au second ordre, l'écart vaut $\Var(\sigma_R^2)/(8V^{3/2})$~: environ 1.6 points à un an
selon un \omterm{def:dv:stochastic-volatility:heston}{modèle de Heston} calibré sur une surface d'indice typique, davantage avec une \omterm{def:dv:stochastic-volatility:sv}{volatilité de la volatilité} plus élevée.
\iqlookfor{le signe, le développement, et le fait qu'il faut un modèle.}
\end{solution}

\begin{solution}{iq:dv:variance-swaps-and-volatility-derivatives:4}
Retrouver le forward par la parité call-put, $K_0$ au niveau ou en dessous de lui, prendre les puts en dessous et les calls au-dessus (les deux moyennés en
$K_0$), pondérer par $\Delta K/K^2$ avec des différences unilatérales aux extrémités, retrancher $(F/K_0-1)^2/T$, interpoler deux
échéances à trente jours. Cas limites~: le forward exactement sur un strike~; les cours acheteurs nuls et l'endroit où couper les ailes~; espacement
irrégulier des strikes~; cotations périmées~; heures d'échéance en minutes.
\iqlookfor{la formule, la correction en $K_0$ et les cas de données.}
\end{solution}

\begin{solution}{iq:dv:variance-swaps-and-volatility-derivatives:5}
La jambe du swap verse $r^2=0.04$ de variance pour la journée, tandis que la couverture rapporte $2(e^r-1-r)=0.0375$~: une perte de l'ordre de
$|r|^3/3$. Si le saut passe sous le strike le plus bas du strip, la couverture cesse aussi de suivre. Les pertes croissent avec le notionnel
et il n'existe aucune correction dynamique~; la protection passe par la valorisation du risque de saut, des caps et l'achat de puts profonds.
\iqlookfor{le terme cubique, la plage du strip et les caps.}
\end{solution}

\begin{solution}{iq:dv:variance-swaps-and-volatility-derivatives:6}
Le contrat à terme vaut $\E[\mathrm{VIX}_T]$ et la volatilité du swap de \omterm{def:dv:rough-volatility-and-forward-variance-models:fwdvar}{variance forward} vaut $\sqrt{\E[\mathrm{VIX}_T^2]}$~; l'inégalité de Jensen
rend la première plus petite, d'un montant fixé par la volatilité de l'indice. Pour la négocier~: acheter la \omterm{def:dv:rough-volatility-and-forward-variance-models:fwdvar}{variance forward}
(un calendar spread de \omterm{def:dv:variance-swaps-and-volatility-derivatives:vs}{swaps de variance}) contre un contrat à terme vendu, qui est long de la volatilité de l'indice,
et couvrir son \omterm{def:dv:greeks-and-the-hedging-pnl:greeks}{vega} avec des options sur indice~; la position profite si la volatilité réalisée de l'indice dépasse ce que
valorise le forward.
\iqlookfor{Jensen, et à quoi le calendar spread est exposé.}
\end{solution}
